% Lösungen Funktionsklassen und Eigenschaften (zusammenbau v0.4, Heft)
\einheitenkopf{Funktionsklassen und Eigenschaften}

\erg{51}{a) Faktoren null setzen \quad b) $x_1 = 2$, $x_2 = -5$ \quad c) $x(x^2 - 9) = 0$: $x_1 = 0$, $x_2 = -3$, $x_3 = 3$}
\erg{52}{a) $S_y(0 | -18)$; $x^2 - 5x + 6 = 0$ ergibt $x = 2$ und $x = 3$; $3$ kommt doppelt vor, also zwei verschiedene Nullstellen \quad b) $S_y(0 | -2)$; $x^2 - x - 2 = 0$ ergibt $x = -1$ und $x = 2$; $-1$ kommt doppelt vor, also zwei verschiedene Nullstellen}
\erg{53}{a) nach unten \quad b) $f(-2) = 2$, $f(-1) = -1$, $f(0) = -2$, $f(1) = -1$, $f(2) = 2$ \quad c) $f(3) = 2{,}5$; $f'(3) < 0$, weil der Graph dort fällt – der positive Funktionswert sagt nichts über die Steigung}

53b)

\wertetabelle[2,-1,-2,-1,2]{x}{f(x)}{-2,-1,0,1,2}\begin{ksys}[xmin=-3,xmax=3,ymin=-3,ymax=3,klein]\funktion{\x^2-2}{f}\end{ksys}

\erg{54}{a) Ecken $(0 | 3)$, $(\pi | -3)$, $(2\pi | 3)$; Grundseite $2\pi$, Höhe $6$ (nicht $3$): $A = \frac{1}{2} \cdot 2\pi \cdot 6 = 6\pi \approx 18{,}85$ \quad b) Ecken $(\pi | 2)$, $(3\pi | -2)$, $(5\pi | 2)$; Grundseite $4\pi$, Höhe $4$ (nicht $2$): $A = \frac{1}{2} \cdot 4\pi \cdot 4 = 8\pi \approx 25{,}13$ \quad c) falsch: auf der x-Achse ist $1$ Längeneinheit $= 2$ cm, auf der y-Achse wegen $f(-1) = 3$ ist $1$ Längeneinheit $= 1$ cm \quad d) richtig: auf der x-Achse ist $1$ Längeneinheit $= 1{,}5$ cm, auf der y-Achse wegen $f(1) = -1{,}5$ ebenfalls $1$ Längeneinheit $= 1{,}5$ cm}
